\documentclass[10pt,journal,compsoc]{IEEEtran}
\usepackage[utf8]{inputenc}
\usepackage{amsmath,amsfonts,amssymb}
\usepackage{graphicx}
\usepackage{cite}
\usepackage{algorithmic}
\usepackage{textcomp}

\begin{document}

\title{A Formal Analysis of the Logarithmic Median Search}

\author{ }

\IEEEtitleabstractindextext{%
\begin{abstract}
This paper presents a rigorous breakdown of the logarithmic time complexity median search algorithm over two sorted arrays, demonstrating its superiority over trivial linear merges.
\end{abstract}
}

\maketitle
\IEEEdisplaynontitleabstractindextext
\IEEEpeerreviewmaketitle

\section{Introduction}

The problem of finding the median of two sorted arrays is a fundamental algorithmic challenge with applications in distributed databases and order statistics. Given two sorted arrays, $A$ of size $m$ and $B$ of size $n$, the trivial solution involves merging the arrays into a single sorted sequence of size $m + n$ and extracting the median in $\mathcal{O}(m + n)$ time. However, in latency-critical systems, a linear time bound is sub-optimal. This paper presents a formal analysis of a binary search partition algorithm that achieves the theoretical lower bound of $\mathcal{O}(\log(\min(m, n)))$ time complexity.

\section{Problem Definition}

Let $A = \langle a_1, a_2, \dots, a_m \rangle$ and $B = \langle b_1, b_2, \dots, b_n \rangle$ be two monotonically non-decreasing sequences. We define the merged sequence $C = A \cup B$ such that $C$ is sorted and $|C| = m + n$. The objective is to identify the median value of $C$ without explicitly constructing $C$.

The median is formally defined as the value separating the higher half from the lower half of the dataset. If $m+n$ is odd, the median is the value at index $(m+n+1)/2$. If $m+n$ is even, it is the arithmetic mean of the values at indices $(m+n)/2$ and $(m+n)/2 + 1$.

\section{The Logarithmic Partition Strategy}

To achieve sub-linear time complexity, we employ a divide-and-conquer strategy utilizing the properties of sorted sequences. The core insight is that finding the median is equivalent to dividing the combined multi-set into two disjoint subsets, $L$ (left) and $R$ (right), satisfying two conditions:
\begin{enumerate}
    \item $|L| = |R|$ (or $|L| = |R| + 1$ for odd lengths)
    \item $\max(L) \leq \min(R)$
\end{enumerate}

We achieve this by making a partition in array $A$ at index $i$ and a corresponding partition in array $B$ at index $j$. 

\subsection{Partition Invariants}

Let the total length be $N = m + n$. The partition index $i$ in array $A$ uniquely determines the partition index $j$ in array $B$ to maintain the size invariant of subset $L$:
\begin{equation}
j = \left\lfloor \frac{m + n + 1}{2} \right\rfloor - i
\end{equation}

By performing a binary search on the smaller array (assume $|A| \leq |B|$ without loss of generality), we can find the optimal partition index $i \in [0, m]$.

\subsection{Validation of Partition}

For a given $i$ and derived $j$, the partition is valid if and only if the cross-boundary conditions hold:
\begin{align}
A[i-1] &\leq B[j] \\
B[j-1] &\leq A[i]
\end{align}

If $A[i-1] > B[j]$, the partition in $A$ is too far to the right, and we must adjust our binary search boundary leftwards. Conversely, if $B[j-1] > A[i]$, the partition is too far to the left.

\section{Complexity Analysis}

\textbf{Theorem 1.} \textit{The proposed partition algorithm determines the median in $\mathcal{O}(\log(\min(m, n)))$ time and $\mathcal{O}(1)$ auxiliary space.}

\textit{Proof.} Let $m = |A|$ and $n = |B|$. By ensuring we only perform the binary search over the smaller array $A$, the search space is restricted to $[0, m]$. In each iteration of the binary search, the search space is halved. Therefore, the maximum number of iterations is strictly bounded by $\lceil \log_2(m) \rceil$. Since all operations within the loop (index arithmetic, boundary condition checks, and median calculation) execute in constant time $\mathcal{O}(1)$, the overall time complexity is $\mathcal{O}(\log(\min(m, n)))$. The space complexity is $\mathcal{O}(1)$ as no dynamic allocations or recursive stack frames are required. $\blacksquare$

\section{Conclusion}

The logarithmic partition method demonstrates optimal asymptotic behavior for the two-array median problem. By leveraging the inherent order of the input sequences and establishing strict mathematical invariants for cross-array partitioning, we eliminate the need for linear merging, yielding a highly performant algorithm suitable for large-scale data processing.

\end{document}
